% Version 0.23: a complete downstream cohomological assembly.
% Four consolidated endpoints VW1--VW4. This fragment introduces no new macros.
\section{Closing the cohomological window assembly}\label{sec:cohomological-window-closure}

\status{The full implication from the two precise source interfaces
\texttt{JointWin} and \texttt{ThmBc} to \texttt{VWinAt} is proved below,
including the actual fields, all finite prime sums, and uniform thresholds.
The two input theorems themselves remain unreconstructed.  With QP22 this
closes the subsequent fixed-pair window construction.  It does not assert
an unconditional cohomological window or a complete RH/GRH proof.}

The usefulness of this assembly is that it removes a collection of
separate field, parity, conductor, scale and cutoff obligations from the
continuation frontier.  All the bounds concern the \emph{same} number
field $E$ specified by \texttt{IsChiField}; no new compatible-form or
coefficient-field witness is manufactured by an isomorphism.

\subsection{The measured groups and the exact three-part sum}

Fix a number field $E$, a rational prime $p$, an algebraic closure and
choices of places over the places of $E$ above $p$.  These are the
source's \texttt{LData}.  Put
\[
 \mathscr S_p(E,j)=\ker\!\left(
 H^1(G_{E,S_p},(\mathbb Q_p/\mathbb Z_p)(-j))
 \longrightarrow\prod_{v\mid p}
 H^1(E_v,(\mathbb Q_p/\mathbb Z_p)(-j))\right),
\]\[
 \ell_p(E,j)=\operatorname{length}_{\mathbb Z_p}\mathscr S_p(E,j).
\]
Here $G_{E,S_p}$ is the Galois group unramified outside the finite places
above $p$, with the usual allowance at infinity.  The lengths are initially
in $\mathbb N\cup\{\infty\}$.  This is exactly the source's negative twist:
\texttt{qzTw D j} uses the cyclotomic character to power $-j$.
We retain all the choices in each argument.  Every bound in this section
is required for \emph{every} algebraic closure and every such system of
local choices, as in \texttt{VVle}; no independence theorem is needed to
replace this quantifier.

For a CM field $E$, complex conjugation acts on $\mathscr S_p(E,j)$.
Write $\ell_p^+(E,j)$ and $\ell_p^-(E,j)$ for the lengths of the
eigenspaces with eigenvalues $(-1)^j$ and $(-1)^{j+1}$, respectively.
Thus the labels $+$ and $-$ include the twist parity; they are not
always the eigenvalues $+1$ and $-1$ of the untwisted action.

For an integer $t\ge1$, use the following notation for the literal source
conditions.  Each condition includes finiteness of every length being
summed, for $1\le j<t$:
\begin{align*}
 \mathcal B_2(E,t;B_2)&:
       \sum_{1\le j<t}\ell_2(E,j)\le B_2,\\
 \mathcal B_\pm(E,t;B_\pm)&:
       \sum_{p\in S}\log p\sum_{1\le j<t}\ell_p^\pm(E,j)\le B_\pm
       \quad\text{for every finite set $S$ of odd primes},\\
 \mathcal B(E,t;B)&:
       \sum_{p\in S}\log p\sum_{1\le j<t}\ell_p(E,j)\le B
       \quad\text{for every finite set $S$ of primes}.
\end{align*}
These are \texttt{BoundedHeightTwoPlaceCondition},
\texttt{VVplusLe}/\texttt{VVminusLe}, and \texttt{VVle}, respectively.
In particular, $\mathcal B_2$ has no factor $\log2$ in its definition.

\begin{proposition}[VW1: exact cohomological assembly]\label{prop:vw-parts}
For every CM field $E$, every integer $t\ge1$, and every specified
system of the preceding data, the three conditions
$\mathcal B_2(E,t;B_2)$, $\mathcal B_-(E,t;B_-)$ and
$\mathcal B_+(E,t;B_+)$ imply
\[
                 \mathcal B(E,t;B_2\log2+B_-+B_+).
                 \tag{VW.1}
\]
This implication preserves all the universal local-data quantifiers.
\end{proposition}

\begin{proof}
Here are the cohomological details needed for the splitting.  Lift an
automorphism $\tau$ of $E$ to an automorphism $\sigma$ of its algebraic
closure, and let $\chi_p$ be the cyclotomic character.  On a continuous
$1$-cocycle $z$ with values in the twist $-j$, the induced map is
\[
 (T_\sigma z)(g)=\chi_p(\sigma)^{-j}z(\sigma^{-1}g\sigma).
 \tag{VW.2}
\]
Conjugation preserves the subgroup generated by inertia outside $p$,
since it permutes places above each rational prime, so this formula
descends to $G_{E,S_p}$.  The cocycle identity shows that it takes
cocycles to cocycles and coboundaries to coboundaries.  If $\sigma$ fixes
$E$, the same identity gives
\[
 T_\sigma z(g)-z(g)=(\rho(g)-1)z(\sigma),
 \qquad \rho=\chi_p^{-j},
\]
which is a coboundary.  Consequently the action on cohomology is
independent of the lift and is multiplicative in $\tau$.  Local
restriction commutes with (VW.2), with the places permuted.  It therefore
preserves the kernel $\mathscr S_p(E,j)$.  For complex conjugation $c$,
this proves $c^2=1$ on that kernel, for the original local data.

When $p$ is odd, $2$ is a unit in $\mathbb Z_p$.  The two endomorphisms
\[
 P_+=\frac{1+(-1)^jc}{2},\qquad
 P_-=\frac{1-(-1)^jc}{2}
\]
are complementary idempotents.  Thus the module itself is the direct
sum of the two eigenspaces, and additivity of module length gives
\[
                \ell_p(E,j)=\ell_p^+(E,j)+\ell_p^-(E,j).
                \tag{VW.3}
\]
This equality holds for extended lengths as well.  The assumed finiteness
of the two summands makes every full odd-primary length finite.
The two-primary finiteness is supplied separately by $\mathcal B_2$.

Now split any finite prime set into its possible member $2$ and its odd
members.  Apply (VW.3) term by term on the odd part, and bound the
$2$-term by $B_2\log2$.  This proves (VW.1), including the finiteness
conjunct in \texttt{VVle}.  The argument applies separately to every
algebraic closure and every local-data system, so none is changed or
selected anew.  No duality or automorphy theorem is used in this assembly.
\end{proof}

\subsection{The actual field and the uniform numerical budget}

Fix a primitive target character $\psi$ of conductor $M>0$.
The condition \texttt{IsChiField}$(\psi,e,q,d;E)$ means that $E$ embeds
in a cyclotomic field and that its fixing automorphisms are exactly those
whose exponents $a$ satisfy
\[
 a\equiv\pm1\pmod{2^{e+2}},\quad \psi(a)=1,\quad
 (a/q)=(a/d)=1.
 \tag{VW.4}
\]
For prime $q\equiv3\pmod8$ and prime $d\equiv7\pmod{16}$ coprime to
$M$, this identifies its image with the compositum
\[
 E_e=\mathbb Q(\zeta_{2^{e+2}})^+
          \mathbb Q(\zeta_M)^{\ker\psi}(\sqrt{-q},\sqrt{-d}).
 \tag{VW.5}
\]
Let $N(E)$ be the least positive $N$ admitting a rational embedding
$E\hookrightarrow\mathbb Q(\zeta_N)$; this is the literal
\texttt{condF E}.  The existence of such $N$ is part of (VW.4).

For $0<\epsilon<1/52$, call $e(x)$ a page function if
\[
 x^{2+8\epsilon}\le2^{e(x)}\le2x^{2+8\epsilon}
                  \quad(x\ge1).
 \tag{VW.6}
\]
The explicit choice
$e(x)=\lceil(2+8\epsilon)\log_2x\rceil$ works for $x\ge1$.
A \emph{window pair at $w$} means the full \texttt{WinPair} condition:
$q,d$ are primes with the congruences just stated,
\begin{gather*}
 \tfrac12w^{\epsilon/5}<\log q\le w^{\epsilon/5},\qquad
 \tfrac12w^{\epsilon/5}(\log w)^2<\log d
                   \le w^{\epsilon/5}(\log w)^2,\\
 q\mid d+1,\qquad
 p\mid d+1\quad\text{for every odd prime }p\le w^{\epsilon/5}.
 \tag{VW.7}
\end{gather*}
The divisibility clauses are retained even though the estimates immediately
below only need the displayed size bounds.

\begin{proposition}[VW2: same-field and uniform numerical lemmas]
\label{prop:vw-budget}
The following assertions hold, with constants depending only on the
parameters declared fixed.
\begin{enumerate}
\item Every field satisfying (VW.4), with $q$ prime and $q\equiv3\pmod4$,
is a CM field and satisfies
\[
 1\le N(E)\le2^{e+2}Mqd,\qquad [E:\mathbb Q]\le N(E).
 \tag{VW.8}
\]
For the host assumptions in (VW.5) and $e\ge2$, its conductor is exactly
\[
                 N(E)=\operatorname{lcm}(2^{e+2},M)qd.
                 \tag{VW.9}
\]
\item Fix $C_t\ge1$.  A page function and $t(x)=\lfloor x^{C_t}\rfloor$
satisfy, for every $x\ge2$,
\[
 e(x)\ge2,\qquad 2^{e(x)}\le x^4,\qquad
             2\le t(x)\le x^{C_t}.
 \tag{VW.10}
\]
\item Uniformly over every window pair at $w$, for sufficiently large
$w$ one has $\log(qd)\le w^{\epsilon/4}$.  For fixed $J'>0,C_t,P$,
all window pairs at $w=X^{1/J'}$, once $X$ is sufficiently large, satisfy
\[
 q,d>2X^{C_t}+1,\qquad q,d>P.
 \tag{VW.11}
\]
\item Uniformly over $E,e,q,d,t$ satisfying (VW.6), (VW.8),
$1\le t\le x^{C_t}$, and $\log(qd)\le x^{\epsilon/4}$, one has
\[
                    \log(N(E)t)\le2x^{\epsilon/4}
                    \tag{VW.12}
\]
for all sufficiently large $x$.
\end{enumerate}
\end{proposition}

\begin{proof}
The fixing condition in (VW.4) and finite Galois correspondence identify
the image of $E$ with the compositum (VW.5), inside the \emph{same}
cyclotomic ambient field.  For the assertion of being CM, it suffices to
see $\sqrt{-q}$ in this image.  With $\zeta_q$ a primitive $q$th root,
the quadratic Gauss sum
$g=\sum_{a\bmod q}(a/q)\zeta_q^a$ satisfies
$\tau_b(g)=(b/q)g$.  Moreover $|g|^2=q$ by additive-character
orthogonality: writing $a=tb$ in the double sum gives $q-1$ for $t=1$
and $-1$ otherwise, and the sum of the nontrivial quadratic character
is zero.  Since $\overline g=(-1/q)g=-g$, it follows that $g^2=-q$.
Every automorphism fixing $E$ fixes $g$ by (VW.4), so $g$ lies in the
image of $E$.  Therefore $E$ has no real embedding.  It is an abelian
subextension of a cyclotomic extension.  Complex conjugation is then a
central nontrivial involution, its fixed field is totally real, and
$E$ is quadratic over that fixed field.  This proves the CM assertion.

Put $L=\operatorname{lcm}(2^{e+2},M,q,d)$.  Every automorphism trivial
modulo $L$ satisfies (VW.4); Galois correspondence gives an embedding
$E\hookrightarrow\mathbb Q(\zeta_L)$.  Hence $N(E)\le L\le2^{e+2}Mqd$.
At the minimal conductor, dimension of the embedding gives
$[E:\mathbb Q]\le\varphi(N(E))\le N(E)$, proving (VW.8).

For (VW.9), the character group of the compositum contains $\psi$, the
primitive quadratic characters of conductors $q,d$, and a primitive
even character of conductor $2^{e+2}$ from the real cyclotomic factor.
Such an even character is obtained by sending $5$ to a primitive
$2^e$th root and $-1$ to $1$ in
$(\mathbb Z/2^{e+2}\mathbb Z)^\times
 =\langle-1\rangle\times\langle5\rangle$.
The identity $v_2(5^{2^r}-1)=r+2$, proved by induction by squaring,
shows that this character does not factor modulo $2^{e+1}$.
Every cyclotomic modulus for the compositum must therefore be divisible
by $M,q,d,2^{e+2}$.  Conversely their least common multiple is a
cyclotomic modulus by the preceding paragraph.  Since $q,d$ are distinct
and prime to $2M$, this proves (VW.9).  This argument is invariant under
the given rational embedding of $E$ and proves the equality for the
original field, not just for a separately chosen model.

For (VW.10), $2+8\epsilon<3$.  At $x\ge2$, (VW.6) gives
$2^{e(x)}\ge x^2\ge4$ and
$2^{e(x)}\le2x^3\le x^4$.
The floor inequalities follow from $C_t\ge1$.

The upper pair bounds give
$\log(qd)\le w^{\epsilon/5}(1+\log^2w)$.
Since $1+\log^2w\le w^{\epsilon/20}$ eventually, the first estimate
in part 3 follows uniformly in the pair.  Set
$\delta=\epsilon/(5J')>0$.  For $w=X^{1/J'}$ and $\log w\ge1$,
both lower pair bounds give $q,d>\exp(X^\delta/2)$.
That expression eventually exceeds $2X^{C_t}+1$ and any fixed $P$.
This proves (VW.11) uniformly over all allowed pairs.

Finally (VW.6) and (VW.8) imply
\[
 \log(N(E)t)\le\log(8M)+(2+8\epsilon+C_t)\log x+\log(qd).
\]
For large $x$, the first two terms are at most $x^{\epsilon/4}$,
independently of $E,e,q,d,t$; the assumed bound on the last term proves
(VW.12).  All ``sufficiently large'' thresholds in this proof involve
only fixed real parameters and $M$, as required.
\end{proof}

\paragraph{The exact odd-minus expression.}
For $o,n,N,t\in\mathbb N$ and $C_b\in\mathbb R$, write
\[
 \mathfrak M(C_b;o,n,N,t)
 =3nt^2\log(tN)+C_b\sum_{\substack{\ell\mid o\\\ell\text{ odd prime}}}
          \ell n^2t^2\log(t\ell N)^2.
 \tag{VW.13}
\]
This is the source's \texttt{minusBound}.  If $1\le N,t$, $n\le N$,
and $\log(Nt)\le2W$, then
\[
 \mathfrak M(C_b;o,n,N,t)\le K\exp(8W),\qquad
 K=3+|C_b|\sum_{\substack{\ell\mid o\\\ell\text{ odd prime}}}\ell^3.
 \tag{VW.14}
\]
Indeed put $u=Nt\ge1$.  The first term is at most $3u^4$ because
$nt\le u$, $t\le u$ and $\log u\le u$.
For the summand at $\ell$, use
$\log(t\ell N)\le\ell u$ and $nt\le u$, obtaining at most
$\ell^3u^4$.  Replace $C_b$ by $|C_b|$ and use $u^4\le e^{8W}$.
The crucial point is that $o$ and $C_b$, hence $K$, are fixed before
$e,q,d,x$ vary.

\paragraph{A uniform final budget.}
Fix real $C$ and $K\ge0$.  If $u\ge1$ and
$\log u\le2x^{\epsilon/4}$, then eventually
\[
 e^{x^{4\epsilon/5}}\log2
       +Ke^{8x^{\epsilon/4}}+u^C+Ce^{x^{4\epsilon/5}}
                     \le e^{x^\epsilon/2}.
 \tag{VW.15}
\]
For a fully explicit threshold set
$A=\log2+|C|+K+1$, $D=9+2|C|$, and require
\[
 x\ge\max\{1,[2(D+\log A+1)]^{5/\epsilon}\}.
\]
Then $b=x^{4\epsilon/5}\ge1$ and $x^{\epsilon/4}\le b$.
Each exponential on the left of (VW.15), and also $u^C$, is bounded
by $e^{Db}$ after replacing coefficients by their absolute values.
The sum is at most $Ae^{Db}$.  The threshold gives
$x^\epsilon/2\ge b(D+\log A+1)\ge Db+\log A$, proving (VW.15).

\subsection{Two precise input interfaces and their closed consequence}

We now state the input interfaces in ordinary mathematical notation.
Their names refer to literal source predicates, rather than additional
assumptions on the final theorem.

Fix $C_t\ge1$.  A pair of functions $(e,t)$ is \emph{admitted with
exponent $A_1$} when, for all $x\ge2$,
\[
 e(x)\ge2,\quad 2^{e(x)}\le x^{A_1},\quad
                   2\le t(x)\le x^{C_t}.
\]
Let $M_\psi=\mathbb Q(\zeta_M)^{\ker\psi}$ and
$o_\psi=[M_\psi:\mathbb Q]$; this equals the order of $\psi$.

\begin{description}[style=nextline,leftmargin=0pt,labelwidth=0pt,labelsep=0pt]
\item[Joint window input $\mathsf J(C_t)$.]
For each primitive $\psi$, each $0<\epsilon<1/52$, each real $A_1$,
and each admitted $(e,t)$, there are constants
$J>0,C,P,X_J$ such that the following holds for every $X\ge X_J$.
Whenever $q,d$ are primes, $q\equiv3\pmod8$, $d\equiv7\pmod{16}$,
$\gcd(qd,M)=1$,
\[
 q,d>2X^{C_t}+1,\quad q,d>P,\quad
                  \log(qd)\le X^{\epsilon/(4J)},
\]
there is $Y$ with $X^{1/J}\le Y$ and $Y^2\le X$ such that,
for every $x\in[Y,Y^2]$ and every field $E$ satisfying
\texttt{IsChiField}$(\psi,e(x),q,d;E)$,
\[
 \begin{split}
 \mathcal B_2(E,t(x);e^{x^{4\epsilon/5}}),\
 \mathcal B_+(E,t(x);(N(E)t(x))^C+Ce^{x^{4\epsilon/5}})
 \end{split}
 \tag{VW.16}
\]
hold with their full local-data quantifiers.  This is
\texttt{JointWin}$(C_t)$: its \texttt{BaAt} and \texttt{HVAt} concern
the same $E$, using
$\texttt{IsCompField}(\texttt{AbF.ofChar}\,\psi)\leftrightarrow
\texttt{IsChiField}\,\psi$.

\item[Odd-minus input $\mathsf M(C_t)$.]
The literal \texttt{ThmBc}$(C_t)$ supplies a constant $C_b$,
uniform over its abelian base fields and admitted functions, such that
for the same fields and parameters
\[
 \mathcal B_-(E,t(x);
       \mathfrak M(C_b;o_\psi,[E:\mathbb Q],N(E),t(x)))
 \tag{VW.17}
\]
holds whenever $x\ge2$ and the prime congruence and coprimality
conditions hold.  Only its specialization to the fixed base $M_\psi$
is needed here; the universal source statement is stronger.
\end{description}

\begin{theorem}[VW3: complete $\mathsf J+\mathsf M$ to $\mathtt{VWinAt}$ chain]
\label{thm:vw-closed}
Let $C_t\ge1$.  Assuming the two interfaces $\mathsf J(C_t)$ and
$\mathsf M(C_t)$ just stated, the literal \texttt{VWinAt}$(\psi,C_t)$
holds for every primitive $\psi$.
More precisely, for every $0<\epsilon<1/52$ and every page function
$e$, there exist $J'\ge1$ and $X_0$ such that for every $X\ge X_0$
and \emph{every} window pair $q,d$ at $w=X^{1/J'}$, there is $Y$ with
\[
 w\le Y,\qquad Y^2\le X,
\]
and for all $x\in[Y,Y^2]$, every field $E$ satisfying
\texttt{IsChiField}$(\psi,e(x),q,d;E)$ satisfies
\[
       \mathcal B(E,\lfloor x^{C_t}\rfloor;e^{x^\epsilon/2}).
       \tag{VW.18}
\]
This includes finiteness of every individual length and bounds every
finite prime subsum, for every algebraic closure and every local-data
system.  The same pair $q,d$ is used throughout the entire square window.
\end{theorem}

\begin{proof}
Fix $\psi,\epsilon,e$.  Put $t(x)=\lfloor x^{C_t}\rfloor$.
By (VW.10), $(e,t)$ is admitted with $A_1=4$.
Apply $\mathsf J(C_t)$ to obtain $J,C,P,X_J$, and let $C_b$ be
provided by $\mathsf M(C_t)$.  Fix
\[
 J'=\max(J,1),\qquad P'=\max(P,M),\qquad
 K=3+|C_b|\sum_{\substack{\ell\mid o_\psi\\\ell\text{ odd prime}}}\ell^3.
\]
These constants are all fixed before $X,q,d,Y,x,E$ are chosen.
Choose $L\ge2$ above the thresholds for the first pair estimate in
VW2, (VW.12), and (VW.15).  Increase $X_0$ to exceed
$X_J,L^{J'}$ and the threshold in (VW.11) for $J',C_t,P'$.

Now let $X\ge X_0$ and let $q,d$ be \emph{any} window pair at
$w=X^{1/J'}$.  Then $w\ge L$ and
\[
 \log(qd)\le w^{\epsilon/4}
       =X^{\epsilon/(4J')}\le X^{\epsilon/(4J)}.
\]
By (VW.11), $q,d>2X^{C_t}+1$ and $q,d>P'$.  Since $M>0$ and these
primes are greater than $M$, they are coprime to $M$.  Thus every
hypothesis of $\mathsf J(C_t)$ holds for this same pair.  It supplies
$Y$ with $X^{1/J}\le Y$, $Y^2\le X$, and (VW.16) on $[Y,Y^2]$.
Because $J'\ge J>0$ and $X\ge1$,
$w=X^{1/J'}\le X^{1/J}\le Y$.

Fix any $x\in[Y,Y^2]$ and any field $E$ in (VW.18).
VW2 proves directly from its fixing condition that this $E$ is CM.
It also gives $1\le N(E)$ and $[E:\mathbb Q]\le N(E)$.
Moreover
\[
 \log(qd)\le w^{\epsilon/4}\le x^{\epsilon/4},\qquad x\ge L.
\]
Hence $u=N(E)t(x)\ge1$ satisfies (VW.12).
The odd-minus input (VW.17), followed by (VW.14) with
$W=x^{\epsilon/4}$, gives
$\mathcal B_-(E,t(x);Ke^{8x^{\epsilon/4}})$.
Use (VW.16) for the other two parts and combine them by VW1.
The resulting upper bound is exactly the left side of (VW.15), so
(VW.18) follows.  No stage chooses a different field, pair, algebraic
closure or local-data system.  Since all thresholds were chosen before
the arbitrary pair, the universal \texttt{WinPair} quantifier is preserved.
\end{proof}

\subsection{The minimal supplier boundary and the fixed-pair endpoint}

The new theorem removes \texttt{VWinAt} as a separate unresolved
\emph{assembly} once $\mathsf J$ and $\mathsf M$ are available.  It does
not construct the cohomological estimates (VW.16)--(VW.17).
For continuation, the source provides a more structural choice of input
boundary:
\[
 \boxed{\quad
 \mathtt{ThmBab}(K,C_t)\ +\ \mathtt{PointHV}\ +\ \mathtt{ThmBcFree}
 \quad}\ .
 \tag{VW.19}
\]
Here \texttt{ThmBab} supplies, on one common square window, the
$2$-primary bound and the finitely many specified odd-primary layer
bounds; \texttt{PointHV} converts those layer bounds into the
odd-plus estimate.  \texttt{ThmBcFree} is precisely the odd-minus
bound (VW.13) for all integer $e,t\ge2$, before admitting scale functions.
These are substantive inputs, not consequences of the arguments above.

\begin{proposition}[VW4: structural supplier and final fixed-pair closure]
\label{prop:vw-fixed-pair}
For $C_t\ge1$, the literal interfaces in (VW.19) imply
$\mathsf J(C_t)$ and $\mathsf M(C_t)$, hence VW3.
For $C_t=401$, combining this with QP22 gives, for every primitive
$\psi$, every $0<\epsilon<1/52$, and every $c\ge1$, a $J_0\ge1$
such that arbitrarily large $Z$ admit a \emph{single} $Y_0,q,d$ with
\[
 \begin{gathered}
 x^{1/J_0}\le Y_0\le x,\qquad
       \mathtt{PairGood}(M,\epsilon,Y_0,x;q,d),\\
 x^{2+8\epsilon}\le c2^{e(x/k)}\le2x^{2+8\epsilon},\qquad
 \mathcal B(E,\lfloor x^{400}\rfloor;e^{x^\epsilon/2})
 \end{gathered}
 \tag{VW.20}
\]
for every $x\in[Z,Z^{3/2}]$ and every
$E$ satisfying \texttt{IsChiField}$(\psi,e(x/k),q,d;E)$, where
$k=c^{1/(2+8\epsilon)}$.  All clauses of \texttt{PairGood}, including
its zero-free boxes and all primitive-inducer clauses, are retained.
When QP22 is used, its named classical inputs CN, BD and HD remain the
analytic leaves of this fixed-pair chain.
\end{proposition}

\begin{proof}
We first check the structural supplier, rather than merely quote an
arrow between predicate names.  For fixed $\psi$, put
$B'=\mathtt{Bchi}\,\psi$ and
$L_\ell=\mathbb Q(\zeta_{\ell M})$ for $\ell\in B'$.
Explicitly, $B'$ consists of $3$ and the odd prime divisors of
$o_\psi M p_Kh_*p_1$.  Here $p_K$ is the least prime at least $7$,
congruent to $3$ modulo $4$ and coprime to $M$; $h_*$ is the class
number of $\mathbb Q(\sqrt{-p_K})$; and $p_1$ is the least prime
congruent to $1$ modulo the odd part of $o_\psi$ and coprime to
$2Mp_K$.  These fixed constants are the ones in GH7 and
\texttt{StmtFields}.
Every $\ell\in B'$ is an odd prime and $L_\ell$ contains $\mu_\ell$.
These are the exact finite set and fields in the source.
The record $K$ specifies nonnegative integers $R_0,R_\ell$ and $C_K>0$.
Put
\[
 R=R_0(M_\psi)+\sum_{\ell\in B'}R_\ell(L_\ell),\qquad
 J_K=2^{R+2}(2C_KC_t/\epsilon)^R,
 \qquad j_\ell=\lfloor\log_\ell R_\ell(L_\ell)\rfloor+3,
 \tag{VW.21}
\]
where the source's natural-number logarithm is $0$ when its argument is
$0$.  The integer $j_\ell$ is therefore defined even when
$R_\ell(L_\ell)=0$.

For the literal \texttt{PointHV} hypothesis, first fix $\psi$ and the
function $\ell\mapsto j_\ell$.  It supplies integers $C,r_1,r_2$ with
$r_1,r_2$ prime.  For every host
$h=(\psi,e,q,d)$ with $e\ge2$ and the prime congruence/coprimality
conditions, it asserts the following.  If $q,d$ exceed $2t+1$,
$\psi.N,\operatorname{ord}(\psi),p_K,h_*,p_1,r_1,r_2$, and if the
specified $\ell$-primary length sum on every prescribed cyclotomic
layer of degree $\ell^{j_\ell}$ over $L_\ell$ is bounded by the same
$B\ge0$ for each $\ell\in B'$, then
\[
       \mathcal B_+(E,t;(N_h t)^C+CB)
       \quad\text{for every corresponding }E,
 \qquad N_h=\operatorname{lcm}(2^{e+2},M)qd.
 \tag{VW.22}
\]
Here the layer quantifier includes every field model $L_0$ of $L_\ell$
and every finite extension $L_1/L_0$ generated by roots of unity of
$\ell$-power order with $[L_1:L_0]=\ell^{j_\ell}$, and retains the
universal algebraic-closure and local-data quantifiers of
$\mathcal B_\ell$ (the unweighted $\ell$-primary length sum).
None of $p_K,h_*,p_1,r_1,r_2$ depends on the varying host.

Apply \texttt{ThmBab} to the abelian base $M_\psi$, this $B'$, and these
$L_\ell$.  It supplies a square window with exponent $J_K$, the
$2$-primary bound in (VW.16), and all the layer bounds just described,
with $B=e^{x^{4\epsilon/5}}$, on the \emph{same} window and for the
\emph{same} pair.  Choose the fixed lower prime cutoff $P$ above the
seven reference numbers in the preceding paragraph.  Increase the
threshold for $X$ so $X^{1/J_K}\ge2$ and the thresholds of both
inputs hold.  Because $t(x)\le x^{C_t}\le X^{C_t}$ on the window,
$q,d>2X^{C_t}+1$ ensures the host inequalities $q,d>2t(x)+1$.
The remaining host inequalities follow from $q,d>P$.
Thus (VW.22) applies at every $x$ in this very window.  The exact
conductor identity (VW.9) gives $N_h=N(E)$ for each corresponding
original field $E$.  Consequently (VW.22) is the odd-plus bound in
(VW.16), with no loss in exponent or conductor.  This proves
$\mathsf J(C_t)$, including its dependence order.  The same field is
used on both sides of the equality; neither a cohomology group nor a
coefficient system has been transported to a newly chosen host.
The free odd-minus input is applied to the integer values $e(x),t(x)$;
(VW.10), or the definition of an admitted pair, proves their lower
bounds.  This gives $\mathsf M(C_t)$ directly.

For (VW.20), use VW3 with $C_t=401$ to fix $J_0$ and its threshold
for the page function.  Apply QP22 with polynomial-window exponent
$J_0$.  Given a lower bound for $Z$, choose $X$ large enough that
$Y_0=X^{1/J_0}$ exceeds that bound, the QP22 selection threshold,
$k^3$ and $k^{400}$.  QP22 gives one pair $q,d$ satisfying the full
\texttt{PairGood} condition on $[Y_0,X]$.  Its \texttt{WinPair} part
allows VW3 to supply $W\ge Y_0$, $W^2\le X$, for this same pair.
Set $Z=kW$.

For $Z\le x\le Z^{3/2}$, the condition $W\ge k^3$ gives
$x\le W^2$, and with $u=x/k$ one has $W\le u\le x\le W^2$.
Restrict the good pair's upper endpoint from $X$ to $x$.
The inequalities $x\le X$ and $Y_0\le W\le x$ give
$x^{1/J_0}\le Y_0\le x$.
Multiplying (VW.6) at $u$ by $c=k^{2+8\epsilon}$ proves the degree
bounds in (VW.20).  Since $u\ge W\ge k^{400}$,
\[
                   x^{400}=k^{400}u^{400}\le u^{401}.
\]
VW3 at $u$ therefore bounds a longer twist range than needed at $x$;
all lengths and prime weights are nonnegative, so the bound persists
on decreasing the cutoff to $\lfloor x^{400}\rfloor$.
Finally $u^\epsilon\le x^\epsilon$ increases the allowed upper bound.
This proves (VW.20) with $Y_0,q,d$ fixed throughout the window.
For ordinary RH, $c=1$ permits retaining the whole square window
$[W,W^2]$, as in OH3.
\end{proof}

\paragraph{The closed chain and its open leaves.}
The chain established in this section is
\[
 \begin{gathered}
 (\mathtt{ThmBab},\mathtt{PointHV},\mathtt{ThmBcFree})
 \Longrightarrow(\mathtt{JointWin},\mathtt{ThmBc})
 \Longrightarrow\mathtt{VWinAt},\\
 \mathtt{VWinAt}+\mathrm{QP22}
 \Longrightarrow\text{one fixed good pair and cohomological window}.
 \end{gathered}
\]
It leaves exactly the deep supplier interfaces on the left (or the
coarser two-interface boundary in VW3).  The source proves them using
its Iwasawa, cohomological counting, congruence and automorphic machinery;
that work has not been reconstructed by this assembly.  In particular,
\texttt{LinkN}/\texttt{LinkV}, the required counting statements and the
actual odd-minus estimate cannot be retired merely because VW3 is now
proved.  Likewise (VW.20) is an input to the later arithmetic geometry;
it does not itself supply the lattice-gain inequality needed to bound
$H_\psi$.
